\section{Considerações Finais e Trabalhos Futuros} \label{sec:conclusion}

Neste trabalho, apresentamos o FLASC,
uma extensão do algoritmo FLSC que introduz um limiar adaptativo $\rho$ para a seleção de modelos e uma fusão local ponderada pela perda empírica.
A motivação central foi proporcionar maior flexibilidade de agrupamento
ao permitir que a quantidade e a influência de cada submodelo global fossem determinadas dinamicamente com base no desempenho local.

Os resultados experimentais conduzidos sobre o conjunto de dados CIFAR-10, em um cenário de heterogeneidade estatística,
revelaram que o FLASC mantém um desempenho equivalente ao FLSC.
Embora a hipótese inicial de que a seleção adaptativa e a ponderação de perdas se traduziria em ganhos significativos,
o FLASC demonstrou que a mudança de abordagem não foi inviável no cenário avaliado.

Como possibilidades para trabalhos futuros, propõem-se as seguintes frentes de investigação:

\begin{itemize}
    \item \textbf{Mais rodadas de treinamento:} Avaliar a convergência do FLASC em simulações de maior duração
    (ex: 20 a 50 rodadas) para verificar se vantagens da fusão adaptativa se revelam em treinamentos prolongados.
    \item \textbf{Exploração extensiva de hiperparâmetros:} Conduzir um estudo que explore mais o parâmetro $\rho$
    para avaliar melhor seu impacto no treinamento.
    \item \textbf{Diversificação do particionamento de dados:} Testar o algoritmo sob outras configurações de heterogeneidade.
\end{itemize}